\chapterimage{images/13-vreme-i-grobnice.jpg}

\chapter{Vreme i Grobnice}

\editorialnote{Razjašnjenja vezana za Vremenske grobnice, Šrajka i prirodu vremena.}

Vremenske grobnice nikada nisu bile samo ruševine.

To je bila prva velika zabluda.

Ljudi su ih vekovima posmatrali kao drevne građevine sačuvane neobičnim poljima.

Kao ostatke neke izgubljene civilizacije.

Kao grobnice u uobičajenom smislu te reči.

Ali antientropijska polja nisu ih samo čuvala od propadanja.

Nosila su ih kroz vreme.

Unazad.

Grobnice nisu dolazile iz prošlosti.

Dolazile su iz budućnosti.

\bigskip

Kasad je to prvi put čuo od Monete još tokom svog ranijeg boravka na Hiperionu.

Tada je mislio da polja sprečavaju starenje građevina.

Moneta ga je ispravila.

``Ne sprečavaju njihovo starenje.''

Pokazala je prema dolini.

``Vraćaju ih kroz vreme.''

Kasad nije odmah razumeo.

Ništa u njegovom vojnom obrazovanju nije ga pripremilo za predmet koji je već postojao zato što je poslat iz vremena koje se još nije dogodilo.

Grobnice su se kretale prema prošlosti sve dok nisu stigle do Hiperiona mnogo pre trenutka u kojem su bile napravljene.

A njihov put još nije bio završen.

\bigskip

Zbog toga su oko njih postojale vremenske plime.

Kada bi polja postala nestabilna, vreme u okolini više nije teklo jednako.

Čovek je mogao da oseti vrtoglavicu.

Da vidi prizore koji još nisu pripadali sadašnjosti.

Da se nađe nekoliko trenutaka, godina ili vekova dalje od mesta na kojem je stajao.

Za većinu ljudi to je bilo samo kratko, dezorijentišuće iskustvo.

Za Rejčel Vejntraub posledice su trajale čitav život.

\bigskip

Kada je kao mlada arheološkinja ušla u Sfingu, temporalni poremećaj zahvatio je njeno telo.

Od tog trenutka njen lični tok vremena krenuo je unazad.

Svaki dan bila je jedan dan mlađa.

Svaki put kada bi zaspala, izgubila bi deo sećanja koji više nije pripadao njenom sve mlađem mozgu.

Ono što je izgledalo kao bolest zapravo je bilo posledica dodira sa mehanizmom koji nije poštovao običan smer vremena.

Rejčel nije samo starila unazad.

Bila je uhvaćena u posledice otvaranja vremenskog prolaza.

\bigskip

Sfinga je bila ključ.

Kada je Severnova proširena svest kasnije mogla da vidi temporalne strukture neposredno, shvatio je da u jednom smislu postoji više Sfingi istovremeno.

Jedna se kretala unazad kroz vreme noseći svoj sadržaj.

Druga je bila nestabilna tačka koja je svojevremeno zahvatila Rejčel.

Treća se već otvorila i počela ponovo da se kreće unapred.

Za ljudsko oko to je bila ista građevina.

Za nekoga ko je mogao da vidi širu strukturu vremena, bila je to ista stvar u različitim delovima sopstvenog putovanja.

\bigskip

I Sfinga je nešto nosila.

Šrajka.

Severn ju je razumeo kao zapečaćeni kontejner koji prenosi smrtonosni sadržaj kroz vekove.

Vremenske grobnice bile su prevozno sredstvo.

Njihov teret nije bio mrtva kultura.

Bio je aktivan.

\bigskip

Šrajk zato nije bio biće koje jednostavno živi na Hiperionu.

Njegova sposobnost da nestaje i pojavljuje se nije bila teleportacija u uobičajenom smislu.

Kretao se kroz vreme.

Kasad je to doživeo neposredno.

Njegovo zaštitno odelo, povezano sa Monetom, uspelo je da prati Šrajka dok se ovaj pomerao između različitih trenutaka.

Jednog časa borili su se u Dolini Vremenskih grobnica.

Sledećeg su bili na vetrokolima na Moru trave.

Za Kasada je promena bila trenutna.

Ali osećaj u kostima i mozgu govorio mu je da su preskočene godine.

\bigskip

Šrajk je mogao da koristi vreme kao čovek prostor.

To je objašnjavalo mnogo toga što je hodočasnicima ranije delovalo nemoguće.

Kako se pojavljivao tamo gde nije mogao da stigne.

Kako je znao gde će ljudi biti.

Kako su događaji razdvojeni danima ili godinama za njega mogli predstavljati susedne trenutke.

Kasad je tokom borbe čak ponovo video prizor sa vetrokola i Heta Mastina.

Shvatio je da ono što je grupa doživljavala kao niz događaja Šrajk možda vidi kao jednu celinu.

\bigskip

To je menjalo i značenje proročanstva.

Ako neprijatelj može da putuje kroz vreme, onda proročanstvo više ne mora da bude predviđanje.

Može biti sećanje.

Poruka poslata iz budućnosti.

Događaj koji je već poznat onome ko ga šalje.

\bigskip

Het Mastin je bio deo takvog proročanstva.

Templari su vekovima verovali da će jedan Pravi Glas Drveta doći na Hiperion.

Dovešće drvobrod.

Videće ga uništenog.

Zatim će uz pomoć erga preuzeti drugo Drvo.

Drvo Ispaštanja.

Drvo trnja.

Šrajkovo Drvo bola.

Mastin je znao da će Igdrasil biti uništen.

Zato njegova mirnoća nakon smrti drvobroda nije bila ravnodušnost.

Bila je ispunjenje nečega što je očekivao.

Ali drugi deo proročanstva nije uspeo.

\bigskip

Kada je video Drvo bola, Mastin nije mogao da prihvati ono što se od njega traži.

Trebalo je da postane njegov Glas.

Da uz pomoć erga vodi tu ogromnu konstrukciju kroz prostor i vreme.

Za Templara, drvobrod je bio živo biće.

Drvo bola bilo je njegova izopačena suprotnost.

Drvo čiji plod nisu bili život i putovanje, već beskrajna ljudska patnja.

Mastin nije bio spreman.

Ili je odbio.

Pobegao je od zadatka za koji je verovao da mu je određen.

\bigskip

Sada je i Drvo bola počelo da ulazi u sadašnjost.

Najpre se pojavljivalo kao nestabilna slika.

Ogromno čelično stablo koje treperi između različitih trenutaka.

Zatim je postajalo sve jasnije.

Na njegovim granama bili su ljudi.

Hiljade.

Desetine hiljada.

Živi.

Svesni.

U bolu.

Martin Silenus je među njima saznao da to nije legenda.

Drvo postoji.

I njegove žrtve ne umiru brzo.

\bigskip

Grobnice su se otvarale.

To nije značilo samo da su se vrata fizički otključavala.

Njihova temporalna izolacija je nestajala.

Ono što je vekovima bilo odvojeno od sadašnjosti ulazilo je u nju.

Drvo bola.

Šrajk.

Prolazi.

Delovi budućnosti.

\bigskip

Bron je, nakon povratka iz megasfere, shvatila da je Drvo bola povezano sa Šrajkovom Palatom.

Nije mogla da objasni kako.

Ali tokom svog prolaska kroz dublje slojeve stvarnosti videla je veze koje običnim čulima nisu postojale.

Palata, Grobnice i Drvo nisu bili samo odvojene građevine razbacane po dolini.

Bili su delovi iste temporalne strukture.

\bigskip

Za Sola je sve to imalo samo jedno značenje.

Rejčel.

Ostalo joj je nekoliko sati.

Zatim minuta.

Kada je Šrajk došao pred Sfingu, Sol je shvatio da se ponavlja san koji ga je godinama progonio.

Ali sada nije bilo oltara iz sna.

Nije bilo glasa iz neba.

Bio je samo on.

Njegova novorođena kćer.

I Šrajk.

\bigskip

Šrajk je otvorio ruke.

Sol je držao Rejčel uz sebe.

Sve što je znao govorilo mu je da ne sme da je preda.

Sve što je osećao govorilo mu je isto.

Ali znao je i da će za nekoliko sekundi njen život završiti ako ništa ne učini.

Njegov izbor više nije bio između sigurnosti i rizika.

Bio je između dve nepoznanice.

\bigskip

Sol se setio sna.

\emph{Kaži da, Tata.}

Godinama je mislio da san od njega zahteva poslušnost.

Sada ga je razumeo drugačije.

Vera nije značila prihvatiti naredbu autoriteta.

Za njega je značila verovati ljubavi čak i kada više nema dokaza da će ona biti nagrađena.

Podigao je Rejčel.

Predao ju je Šrajku.

\bigskip

Šrajk ju je uzeo.

Povukao se u svetlost Sfinge.

I nestao.

Drvo bola iza njih došlo je još jasnije u sadašnjost.

Grobnice su se otvarale.

\bigskip

Ali Sol nije mogao da vidi ono što se dešavalo unutar Sfinge.

Severn jeste.

Do tada Džozef Severn više nije imao telo.

Njegov put sa Lijem Hantom završio se na mestu za koje je čovečanstvo vekovima verovalo da više ne postoji.

Staroj Zemlji.

TehnoCentar je nije uništio u Velikoj Grešci.

Premestio ju je.

Sakrio.

A Severn i Hant našli su se u Rimu iz vremena koje je ličilo na poslednje dane pravog Džona Kitsa.

\bigskip

Severn je tamo ponovo umirao.

Njegovo kibridno telo nije moglo da izdrži.

Hant je ostao uz njega.

Kada je Severn umro, Li Hant ga je odvezao do starog groblja izvan Rima.

Šrajk ih je pratio.

Nije napao.

Samo je išao iza pogrebne kočije.

Hant je sam sahranio telo čoveka kojeg je poznavao kao Džozefa Severna.

Ali smrt tela nije značila kraj ličnosti.

\bigskip

Severn je bio više od kibrida.

Njegov ljudski um bio je povezan sa strukturama TehnoCentra i sa dubljim slojem stvarnosti koji su AI nazivale Prazninom Koja Vezuje.

Kada je telo umrlo, njegova svest nije nestala.

Ostala je bez fizičkog oblika.

Mogla je da prolazi kroz datasferu, megasferu i metasferu.

Da vidi veze koje ljudska čula nisu mogla da vide.

Da prati Hiperion čak i kada se Mreža raspadala.

\bigskip

Zato je, dok je Sol čekao pred Sfingom, Severn mogao da vidi ono što se događalo unutra.

Njegova svest više nije bila ograničena na jedno telo ni na običan tok vremena.

Video je Šrajka kako nosi Rejčel kroz portal.

Tri njegova koraka bila su dovoljna da prođu sati.

Što je dublje ulazio, tok vremena bio je sve snažniji.

Ako nastavi, odneće dete daleko u budućnost.

\bigskip

Severn tada više nije bio običan posmatrač.

Uz pomoć erga mogao je kratko da dobije dovoljno prisustva da deluje.

Prišao je Šrajku.

Uzeo Rejčel iz njegovih ruku.

Šrajk se okrenuo.

Ali bio je preblizu toku vremena koji ga je nosio.

Temporalna struja povukla ga je dalje.

Za nekoliko trenutaka postao je tačka.

Zatim ništa.

\bigskip

Severn je ostao sa novorođenčetom u rukama.

Nije mogao da se vrati istim putem.

Struja vremena bila je previše jaka.

Erg je mogao da pomogne njemu.

Ne i njemu zajedno sa detetom.

Zato je čekao.

Negde između trenutaka.

Sa Rejčel u naručju.

\bigskip

Sol napolju nije znao ništa od toga.

Za njega je Šrajk odneo njegovu kćer.

Seo je pred Sfingu i čekao.

Rat na nebu više mu nije bio važan.

Pad carstva nije mu bio važan.

Sve što je želeo bilo je da se Rejčel vrati.

\bigskip

Sol napolju nije znao ništa od onoga što se dogodilo unutar Sfinge.

Za njega je Šrajk odneo njegovu kćer.

Seo je pred ulaz.

I čekao.

Satima.

Nije ga više zanimao rat.

Ni pad Hegemonije.

Ni TehnoCentar.

Samo Rejčel.

\bigskip

Pred jutro se svetlost u Sfingi promenila.

Sol je ustao.

Neko je izlazio iz portala.

Prvo je video žensku priliku.

Zatim ono što je nosila.

Bebu.

Sol je prestao da diše.

\bigskip

Žena je izašla iz svetlosti.

Imala je oko dvadeset pet godina.

Bakarnosmeđu kosu.

Zelene oči sa sitnim smeđim tačkama.

Lice koje je Sol poznavao bolje od sopstvenog.

``Rejčel.''

Mlada žena se nasmešila.

``Tata.''

U rukama je nosila novorođenče.

Sebe.

\bigskip

Rejčel je nosila Rejčel.

Odrasla žena prišla je ocu i predala mu bebu.

Sol ju je privio uz grudi.

``Da li je...''

Nije uspeo da dovrši pitanje.

``Dobro je'', rekla je odrasla Rejčel. ``Sada normalno stari.''

Sol je pogledao bebu.

Pa odraslu kćer.

Pa ponovo bebu.

Dve verzije istog čoveka stajale su pred njim.

Jedna tek rođena.

Druga odrasla u budućnosti.

\bigskip

U međuvremenu su stigli i ostali.

Bron je pomagala povređenom Martinu.

Konzul i Teo bili su sa njima.

Melio Arundez zastao je kao ukopan kada je video odraslu Rejčel.

Ali Bron ju je prepoznala iz drugog razloga.

Zurila je u mladu ženu.

U kratku kosu.

Zelene oči.

Lice koje je već videla kraj Kasadovog mrtvog tela u Kristalnom Monolitu.

``Moneta.''

Odrasla Rejčel ju je pogledala.

Bron je pokazala prema njoj.

``Ti si Moneta.''

Rejčel je klimnula.

\bigskip

Odgovor koji je Kasad tražio čitavog života konačno je postao jasan.

Moneta nije bila Šrajk.

Nije bila demonski dvojnik žene koju je voleo.

Bila je Rejčel Vejntraub.

Solova kćer.

\bigskip

Rejčel je objasnila da je odrasla u dalekoj budućnosti.

Tamo se vodio konačni rat između čovečanstva i Ultimativne Inteligencije koju je stvorio TehnoCentar.

Njena Merlinova bolest nije bila slučajna nesreća.

Učinila ju je sposobnom za putovanje koje je trebalo da obavi.

Bila je izabrana da se zajedno sa Vremenskim grobnicama vraća kroz vreme.

Da nadzire Šrajka.

I da jednog dana pronađe Fedmana Kasada.

\bigskip

``Ali Kasad te je već poznavao'', rekao je Martin.

Rejčel je odmahnula glavom.

``Za njega jesam.''

Pogledala je prema Kristalnom Monolitu.

``Za mene još nisam.''

Videla je Kasada kako gine u poslednjoj bici.

Ispratila je njegovo telo u Grobnicu.

Ali čoveka koji će je voleti kroz njegove snove i ratove tek je trebalo da upozna.

\bigskip

Kasadova prošlost bila je Rejčelina budućnost.

Rejčelina prošlost tek je trebalo da postane Kasadova prošlost.

Kada ga bude srela, on će je nazvati Monetom.

Ime joj još nije pripadalo.

Tek će ga prihvatiti.

\bigskip

Time se zatvarala vremenska petlja.

Rejčel je kao mlada arheološkinja došla na Hiperion i bila pogođena temporalnim poljem Sfinge.

Dvadeset šest godina starila je unazad.

Kao novorođenče, Sol ju je predao Šrajku.

Severn ju je oteo od Šrajka unutar temporalnog prolaza.

U budućnosti je ponovo odrasla.

Postala je žena koju će Kasad poznavati kao Monetu.

Zatim je sa Vremenskim grobnicama krenula unazad kroz vreme.

Sve do trenutka u kojem je njen sopstveni otac ponovo čeka pred Sfingom.

\bigskip

Sol je jedva uspevao da prati objašnjenje.

``Znači ova beba...''

Pogledao je dete u rukama.

``...će postati ti?''

``Da.''

``A ti sada ideš nazad?''

``Napred za mene. Nazad za vas.''

Sol je zatvorio oči.

Čak ni posle svega nije postojao jezik dovoljno jednostavan za to.

\bigskip

Odrasla Rejčel mu je rekla da ne može dugo da ostane.

Njene dve verzije ne mogu trajno postojati zajedno u istom vremenu.

Morala je da se vrati.

Sol ju je uhvatio za ruku.

``Ne.''

``Tata.''

``Tek sam te dobio nazad.''

Rejčel je dodirnula bebinu glavu.

``I dalje me imaš.''

\bigskip

Objasnila mu je da Sfinga sada može da posluži kao jednosmerni prolaz prema njenom vremenu.

Ako Sol želi, može da pođe sa bebom.

Ali to će značiti da Rejčel mora da podigne još jednom.

Da ponovo proživi njeno detinjstvo.

Treći put.

``Nijedan roditelj ne bi trebalo to da mora'', rekla je.

Sol se nasmešio kroz suze.

``Nijedan roditelj to ne bi odbio.''

\bigskip

Rejčel se vratila prema svetlosti.

Pre nego što je nestala, još jednom je pogledala oca.

``Vidimo se tamo.''

Zatim je otišla.

Sol je ostao sa novorođenom Rejčel u rukama.

Ovoga puta dete je starilo u pravom smeru.

\bigskip

Kasad je na drugom kraju te iste vremenske petlje tek trebalo da razume ko je žena koju voli.

Kada ju je u dalekoj budućnosti pitao za ime, ona ga još nije imala.

``Moneta'', rekao joj je.

Svidelo joj se.

Za njega je njihov odnos već trajao čitav život.

Za nju je tek počinjao.

\bigskip

Zato Kasadova ljubavna priča nije bila pravolinijska.

Njih dvoje su se susretali krećući se kroz vreme u suprotnim smerovima.

On je upoznavao njenu budućnost kao svoju prošlost.

Ona je upoznavala njegovu prošlost kao svoju budućnost.

A oba puta vodila su prema poslednjoj bici.

Moneta koju je poznavao čitavog života nije ga poznavala.

Za njega je njihova ljubav bila prošlost.

Za nju još budućnost.

``Moneta'', rekla je kada je čula ime.

``Dobro ime.''

Kasad ju je gledao.

``To nije tvoje ime?''

``Još nije.''

Tada je konačno razumeo.

Žena koja ga je pratila kroz njegove snove i bitke tek će postati Moneta.

Njeno putovanje kroz vreme bilo je suprotno njegovom.

Njihova veza nije počinjala na istom mestu za oboje.

\bigskip

Ona ga je odvela u daleku verziju Doline Vremenskih grobnica.

Tamo je video hiljade Šrajkova.

I hiljade ljudi spremnih da im se suprotstave.

Nije znao koliko daleko u budućnosti se nalazi.

Moneta mu je rekla samo ono najvažnije.

Rat neće biti ograničen na Hiperion.

Protezaće se preko miliona svetova.

I kroz vreme.

\bigskip

Pobednik neće dobiti samo teritoriju.

Dobiće mogućnost da utiče na prošlost i budućnost.

Ako pobede sile povezane sa Šrajkom, Šrajk koji je već poslat u prošlost biće samo prvi.

Ako pobedi druga strana, čovečanstvo možda dobije priliku da samo odlučuje o sopstvenoj istoriji.

Kasad nije razumeo fiziku.

Nije razumeo metafiziku.

Razumeo je izbor.

Okrenuo se prema Šrajkovima.

I krenuo u borbu.

\bigskip

Tako je postajalo jasno da su Vremenske grobnice mnogo više od portala.

One su bile čvorište.

Mesto na kojem su buduće sile pokušavale da menjaju istoriju.

Jedna strana slala je Šrajka unazad.

Druga je pokušavala da ga zaustavi.

A ljudi u sadašnjosti bili su uhvaćeni između odluka koje još nisu ni donesene.

\bigskip

To je bio najveći paradoks Hiperiona.

Budućnost nije bila posledica sadašnjosti.

Bar ne samo to.

Budućnost je već delovala na sadašnjost.

Slala je bića.

Građevine.

Proročanstva.

Ljubavnike.

Ratnike.

Čitave sudbine.

\bigskip

Zato TehnoCentar nije mogao pouzdano da izračuna Hiperion.

Njegovi modeli počivali su na uzroku koji prethodi posledici.

Na Hiperionu posledica je ponekad već postojala pre uzroka.

Šrajk je bio tamo zato što će jednog dana biti poslat.

Moneta je poznavala Kasada zato što će ga jednog dana upoznati.

Rejčel je obolela zbog portala koji još nije bio potpuno otvoren.

Templari su imali proročanstvo o događajima koji su možda već bili viđeni iz budućnosti.

\bigskip

Hiperion nije kršio samo ljudsko shvatanje vremena.

Krčio je prostor u kojem predviđanje prestaje da bude pouzdano.

Upravo zato je predstavljao Varijablu.

\bigskip

Ali ni Šrajk nije bio potpuno slobodan.

Kasadova borba pokazala je da može biti povređen.

Bron će kasnije pokazati da može biti čak i zaustavljen.

Temporalna moć nije značila svemoć.

Šrajk je bio strašan zato što se kretao pravilima koja ljudi još nisu razumeli.

Ne zato što pravila nisu postojala.

\bigskip

A kada su Grobnice konačno počele da se otvaraju, ta pravila su postajala vidljivija.

Temporalna polja su slabila.

Granice među različitim trenucima postajale su porozne.

Drvo bola ulazilo je u fazu sa sadašnjošću.

Sfinga se pretvarala u aktivan portal.

Ljudi koji su bili razdvojeni godinama mogli su postati gotovo susedi.

\bigskip

Hiperion je postajao mesto na kojem vreme više nije bilo reka koja teče samo u jednom smeru.

Bilo je more.

Sa strujama.

Virovima.

Povratnim tokovima.

I bićima koja su naučila da plove njime.

\bigskip

Za hodočasnike to saznanje nije donosilo utehu.

Naprotiv.

Ako se budućnost može vratiti i menjati prošlost, onda više nije bilo jasno šta znači sudbina.

Da li su njihovi izbori zaista njihovi?

Da li su došli na Hiperion zato što su tako odlučili?

Ili zato što je neko u budućnosti već znao da će doći?

\bigskip

Sol je ipak imao svoj odgovor.

Nije znao da li je Rejčel spasena.

Nije znao šta Šrajk predstavlja.

Nije znao ko je napravio Grobnice.

Ali znao je da je odluku da preda dete doneo on.

Možda je budućnost čekala taj izbor.

Možda ga je čak izazvala.

Ali nije mogla da ga učini umesto njega.

\bigskip

Isto je važilo za Kasada.

Za Bron.

Za Silenusa.

Za Mastina.

Čak i ako su bili deo plana koji se prostirao vekovima, način na koji će odgovoriti tom planu još nije bio potpuno određen.

Tu je ležala slabost svake sile koja pokušava da upravlja istorijom.

Ljudi nisu uvek radili ono što je model predviđao.

\bigskip

Vremenske grobnice bile su mašina budućnosti.

Šrajk njeno oružje.

Ali Hiperion je ostao nepoznanica zato što su se oko njih nalazili ljudi.

I njihove odluke.

A sada, dok su se Grobnice otvarale i budućnost ulazila u sadašnjost, više nije bilo moguće držati ta dva vremena odvojeno.